% Each small diagram shows one numerical elimination update. Values and backtracking rules are in the adjacent text.
\begin{tikzpicture}[x=1mm,y=1mm,>=Stealth,font=\figurefont\fontsize{8}{10}\selectfont,
 v/.style={circle,draw=FigureStroke,fill=FigureBlueFill,line width=.8pt,minimum size=5.5mm,inner sep=0pt},
 cur/.style={v,draw=FigureYellow,fill=FigureYellowFill},
 f/.style={rectangle,draw=FigureStroke,fill=white,line width=.8pt,minimum width=7mm,minimum height=5mm,inner sep=.7pt},
 fn/.style={f,draw=FigureGreen,fill=FigureGreenFill},
 e/.style={draw=FigureStroke,line width=.8pt}]
\begin{scope}[xshift=0mm]
\node[v] (v2) at (9,0) {$2$};
\node[v] (v3) at (18,0) {$3$};
\node[cur] (v4) at (27,0) {$4$};
\node[f] (p1) at (10,14) {$\psi_1$};
\draw[e] (p1)--(v2);
\draw[e] (p1)--(v3);
\node[f] (B) at (23,14) {$\phi_B$};
\draw[e] (B)--(v2);
\draw[e] (B)--(v3);
\draw[e] (B)--(v4);
\end{scope}
\draw[->,FigureMuted,line width=1pt] (33,7)--(40,7);
\begin{scope}[xshift=47mm]
\node[v] (v2) at (9,0) {$2$};
\node[v] (v3) at (18,0) {$3$};
\node[f] (p1) at (7,14) {$\psi_1$};
\draw[e] (p1)--(v2);
\draw[e] (p1)--(v3);
\node[fn] (p4) at (21,14) {$\psi_4$};
\draw[e] (p4)--(v2);
\draw[e] (p4)--(v3);
\end{scope}
\end{tikzpicture}
